\documentclass[11pt]{article}
\usepackage[T1]{fontenc}
\usepackage[a4paper,margin=1in]{geometry}
\usepackage{amsmath,amssymb}
\usepackage[hidelinks]{hyperref}
\usepackage{microtype}

\setlength{\parindent}{0pt}
\setlength{\parskip}{0.65em}

\title{Literature Status: Eigenvalue Convergence for\\
Multiplicative Brownian Motion}
\author{Run \texttt{fa\_banach\_001} --- lane 14}
\date{11 August 2026}

\begin{document}
\maketitle

\begin{center}
\fbox{\parbox{0.9\textwidth}{
\textbf{Status: literature already answered.}
Ho--Zhong ask whether the empirical eigenvalue distribution of
$U_NG_N(t)$ converges to the Brown measure of $ub_t$, including the case
$U_N=I$.  Theorem~1.2 and Definition~1.7 of
Brailovskaya--Cook--Kemp--Parraud prove almost-sure weak convergence to
exactly that Brown measure, in a more general initial-condition setting.
}}
\end{center}

\section{Original question}

Ching-Wei Ho and Ping Zhong, \emph{Brown Measures of Free Circular and
Multiplicative Brownian Motions with Self-Adjoint and Unitary Initial
Conditions}, arXiv:1908.08150, later published in \emph{JEMS}
\textbf{25} (2023), 2163--2227~\cite{source}, compute the Brown measure of
$ub_t$.

On arXiv PDF page~7 they state that it remains open, even for identity
initial data, to prove
\[
  \mu_{G_N(t)} \Longrightarrow \mu_{b_t},
\]
where the left side is the empirical eigenvalue distribution and the right
side is the Brown measure.  On PDF page~50 they state the general unitary
version: for a deterministic, or independent random, unitary $U_N$ whose
empirical eigenvalue distribution tends weakly to the spectral law of a
unitary $u$, prove
\[
  \mu_{U_NG_N(t)} \Longrightarrow \mu_{ub_t}.
\]

\section{Separate later answer}

Tatiana I. Brailovskaya, Nicholas A. Cook, Todd Kemp, and F\'elix Parraud,
\emph{Eigenvalues of Brownian Motions on $\mathrm{GL}(N,\mathbb C)$},
arXiv:2511.10535v2~\cite{answerpaper}, prove the following on arXiv PDF
page~9 (Theorem~1.2).  Let $B(t)$ be their Brownian motion on
$\mathrm{GL}(N,\mathbb C)$ and let $B_0$ be independent of it.  If $B_0$
converges almost surely in $*$-distribution to $b_0$, then, for every fixed
$t>0$, the empirical eigenvalue law of $B_0B(t)$ converges weakly almost
surely to a deterministic measure $\mu_{b_0,t}$.  They explicitly note that
for normal $B_0$, the hypothesis is simply weak convergence of its empirical
eigenvalue distribution.

Definition~1.7 on PDF page~15 identifies $\mu_{b_0,t}$ as the Brown measure
of $b_0b(t)$ and restates Theorem~1.2 in those terms.  Thus $B_0=U_N$ gives
the full unitary-initial statement requested by Ho--Zhong, and $U_N=I$ gives
the identity case.

\section{Provenance and scope}

This is an explicit literature closure, not a new theorem inferred by the
present run.  On PDF page~17, the supporting paper says that Theorem~1.2
solves the convergence problem, identifies the limiting object with the
Brown measures computed in earlier work, and specifically credits
Ho--Zhong as reference~[68] for the unitary-initial Brown-measure computation.

The later theorem is strictly stronger in its initial-data scope: it allows
arbitrary, possibly nonnormal, independent $B_0$ having an almost-sure
$*$-distribution limit.  No part of the source's stated identity or unitary
initial-condition convergence question remains open after this theorem.
A bounded local/arXiv search used the exact paper title, arXiv id, and the
terms ``empirical eigenvalue distribution'', ``multiplicative Brownian
motion'', and ``Brown measure''; arXiv:2511.10535 was the decisive match.

\begin{thebibliography}{9}
\bibitem{source}
C.-W. Ho and P. Zhong,
\emph{Brown measures of free circular and multiplicative Brownian motions
with self-adjoint and unitary initial conditions},
J. Eur. Math. Soc. \textbf{25} (2023), 2163--2227;
arXiv:1908.08150.

\bibitem{answerpaper}
T. I. Brailovskaya, N. A. Cook, T. Kemp, and F. Parraud,
\emph{Eigenvalues of Brownian Motions on $\mathrm{GL}(N,\mathbb C)$},
arXiv:2511.10535v2 (2026).
\end{thebibliography}

\end{document}
